% Options for packages loaded elsewhere
\PassOptionsToPackage{unicode}{hyperref}
\PassOptionsToPackage{hyphens}{url}
\PassOptionsToPackage{dvipsnames,svgnames,x11names}{xcolor}
%
\documentclass[
  11pt,
  letterpaper,
]{article}
\usepackage{amsmath,amssymb}
\usepackage{iftex}
\ifPDFTeX
  \usepackage[T1]{fontenc}
  \usepackage[utf8]{inputenc}
  \usepackage{textcomp} % provide euro and other symbols
\else % if luatex or xetex
  \usepackage{unicode-math} % this also loads fontspec
  \defaultfontfeatures{Scale=MatchLowercase}
  \defaultfontfeatures[\rmfamily]{Ligatures=TeX,Scale=1}
\fi
\usepackage{lmodern}
\ifPDFTeX\else
  % xetex/luatex font selection
\fi
% Use upquote if available, for straight quotes in verbatim environments
\IfFileExists{upquote.sty}{\usepackage{upquote}}{}
\IfFileExists{microtype.sty}{% use microtype if available
  \usepackage[]{microtype}
  \UseMicrotypeSet[protrusion]{basicmath} % disable protrusion for tt fonts
}{}
\makeatletter
\@ifundefined{KOMAClassName}{% if non-KOMA class
  \IfFileExists{parskip.sty}{%
    \usepackage{parskip}
  }{% else
    \setlength{\parindent}{0pt}
    \setlength{\parskip}{6pt plus 2pt minus 1pt}}
}{% if KOMA class
  \KOMAoptions{parskip=half}}
\makeatother
\usepackage{xcolor}
\usepackage[margin=1in]{geometry}
\usepackage{color}
\usepackage{fancyvrb}
\newcommand{\VerbBar}{|}
\newcommand{\VERB}{\Verb[commandchars=\\\{\}]}
\DefineVerbatimEnvironment{Highlighting}{Verbatim}{commandchars=\\\{\}}
% Add ',fontsize=\small' for more characters per line
\newenvironment{Shaded}{}{}
\newcommand{\AlertTok}[1]{\textcolor[rgb]{1.00,0.00,0.00}{\textbf{#1}}}
\newcommand{\AnnotationTok}[1]{\textcolor[rgb]{0.38,0.63,0.69}{\textbf{\textit{#1}}}}
\newcommand{\AttributeTok}[1]{\textcolor[rgb]{0.49,0.56,0.16}{#1}}
\newcommand{\BaseNTok}[1]{\textcolor[rgb]{0.25,0.63,0.44}{#1}}
\newcommand{\BuiltInTok}[1]{\textcolor[rgb]{0.00,0.50,0.00}{#1}}
\newcommand{\CharTok}[1]{\textcolor[rgb]{0.25,0.44,0.63}{#1}}
\newcommand{\CommentTok}[1]{\textcolor[rgb]{0.38,0.63,0.69}{\textit{#1}}}
\newcommand{\CommentVarTok}[1]{\textcolor[rgb]{0.38,0.63,0.69}{\textbf{\textit{#1}}}}
\newcommand{\ConstantTok}[1]{\textcolor[rgb]{0.53,0.00,0.00}{#1}}
\newcommand{\ControlFlowTok}[1]{\textcolor[rgb]{0.00,0.44,0.13}{\textbf{#1}}}
\newcommand{\DataTypeTok}[1]{\textcolor[rgb]{0.56,0.13,0.00}{#1}}
\newcommand{\DecValTok}[1]{\textcolor[rgb]{0.25,0.63,0.44}{#1}}
\newcommand{\DocumentationTok}[1]{\textcolor[rgb]{0.73,0.13,0.13}{\textit{#1}}}
\newcommand{\ErrorTok}[1]{\textcolor[rgb]{1.00,0.00,0.00}{\textbf{#1}}}
\newcommand{\ExtensionTok}[1]{#1}
\newcommand{\FloatTok}[1]{\textcolor[rgb]{0.25,0.63,0.44}{#1}}
\newcommand{\FunctionTok}[1]{\textcolor[rgb]{0.02,0.16,0.49}{#1}}
\newcommand{\ImportTok}[1]{\textcolor[rgb]{0.00,0.50,0.00}{\textbf{#1}}}
\newcommand{\InformationTok}[1]{\textcolor[rgb]{0.38,0.63,0.69}{\textbf{\textit{#1}}}}
\newcommand{\KeywordTok}[1]{\textcolor[rgb]{0.00,0.44,0.13}{\textbf{#1}}}
\newcommand{\NormalTok}[1]{#1}
\newcommand{\OperatorTok}[1]{\textcolor[rgb]{0.40,0.40,0.40}{#1}}
\newcommand{\OtherTok}[1]{\textcolor[rgb]{0.00,0.44,0.13}{#1}}
\newcommand{\PreprocessorTok}[1]{\textcolor[rgb]{0.74,0.48,0.00}{#1}}
\newcommand{\RegionMarkerTok}[1]{#1}
\newcommand{\SpecialCharTok}[1]{\textcolor[rgb]{0.25,0.44,0.63}{#1}}
\newcommand{\SpecialStringTok}[1]{\textcolor[rgb]{0.73,0.40,0.53}{#1}}
\newcommand{\StringTok}[1]{\textcolor[rgb]{0.25,0.44,0.63}{#1}}
\newcommand{\VariableTok}[1]{\textcolor[rgb]{0.10,0.09,0.49}{#1}}
\newcommand{\VerbatimStringTok}[1]{\textcolor[rgb]{0.25,0.44,0.63}{#1}}
\newcommand{\WarningTok}[1]{\textcolor[rgb]{0.38,0.63,0.69}{\textbf{\textit{#1}}}}
\usepackage{longtable,booktabs,array}
\usepackage{calc} % for calculating minipage widths
% Correct order of tables after \paragraph or \subparagraph
\usepackage{etoolbox}
\makeatletter
\patchcmd\longtable{\par}{\if@noskipsec\mbox{}\fi\par}{}{}
\makeatother
% Allow footnotes in longtable head/foot
\IfFileExists{footnotehyper.sty}{\usepackage{footnotehyper}}{\usepackage{footnote}}
\makesavenoteenv{longtable}
\setlength{\emergencystretch}{3em} % prevent overfull lines
\providecommand{\tightlist}{%
  \setlength{\itemsep}{0pt}\setlength{\parskip}{0pt}}
\setcounter{secnumdepth}{-\maxdimen} % remove section numbering
\usepackage{microtype}
\usepackage{etoolbox}
\usepackage{xurl}
\usepackage{needspace}
\setlength{\emergencystretch}{3em}
\AtBeginEnvironment{longtable}{\small}
\pretocmd{\section}{\Needspace{9\baselineskip}}{}{}
\widowpenalty=10000
\clubpenalty=10000
\ifLuaTeX
  \usepackage{selnolig}  % disable illegal ligatures
\fi
\usepackage{bookmark}
\IfFileExists{xurl.sty}{\usepackage{xurl}}{} % add URL line breaks if available
\urlstyle{same}
\hypersetup{
  pdftitle={An Improved Exponent Bound for the Exact Discrete Fourier Transform},
  colorlinks=true,
  linkcolor={Maroon},
  filecolor={Maroon},
  citecolor={Blue},
  urlcolor={Blue},
  pdfcreator={LaTeX via pandoc}}

\title{An Improved Exponent Bound\\for the Exact Discrete Fourier Transform}
\usepackage{etoolbox}
\makeatletter
\providecommand{\subtitle}[1]{% add subtitle to \maketitle
  \apptocmd{\@title}{\par {\large #1 \par}}{}{}
}
\makeatother
\subtitle{Research draft\par\vspace{-22pt}}
\author{}
\date{October 8, 2026}

\begin{document}
\maketitle
\begin{abstract}
The candidate conclusion, in precisely the upstream
arithmetic/root/address model, is

\[
\boxed{T(n)=O\!\left(n(\log n)^{1-5.5\times10^{-10}}\right).}
\]

There is a simpler fallback with exponent saving \(5.2\times10^{-10}\)
that does \textbf{not} depend on the shared-sum replacement or the new
batching lemma. A middle version obtains \(5.3\times10^{-10}\) without
the shared-sum replacement.
\end{abstract}

\section{1. Source, attribution, and the exact
target}\label{source-attribution-and-the-exact-target}

Write {[}O{]} for the uploaded 31-page \emph{An explicit power saving
for the exact discrete Fourier transform}. Its main quantitative proof
is Sections 2--5. The uploaded 49-page companion, \emph{Finite tensor
savings and exact Fourier circuits}, gives an independent finite-win
existence route; that argument is not required for the explicit
quantitative modification here. {[}O{]}, Introduction and Section 1.2,
explicitly separates these routes.

The model counts exact complex field operations, preparation of
input-independent coefficients, schedule construction, array
organization, and unit-cost operations on logarithmic-size integer
addresses. A specified root of unity is supplied as in {[}O{]}. No
bounded-coefficient, stable floating-point, or integer bit-complexity
conclusion is asserted.

{[}O{]} fixes \(h=100\) and proves

\[
T(n)=O\!\left(n(\log n)^{\theta_0}(\log\log n)^{4-\theta_0}\right),
\qquad 1-\theta_0=2.10643843004018\ldots\times10^{-13}.
\]

Its headline \(10^{-13}\) saving is a weaker rounded corollary.
Comparisons must use both quantities honestly.

Douglas Colkitt already identified the usefulness of the unchanged
complex network at \(h=25\), and gave its parameters and compatibility
argument in \texttt{notes/independent-complex.tex} in
\texttt{CrocSwap/integer-mult-bounds}. We do \textbf{not} claim the
independent-ground-size idea or that parameter choice as new. His
displayed complex construction still has \(h+1\) central wires. The
present draft removes a redundant central wire, supplies a no-padding
recurrence, and completes a proposed phase-label integration for the
earlier shared-sum circuit. Originality beyond the sources checked has
not been established.

\section{2. Remove one redundant central
wire}\label{remove-one-redundant-central-wire}

Let

\[
\mathcal T=\binom{[h]}3,\quad v=\binom h3,\quad m=h^3,\quad N=v^3,\quad I=3v^2.
\]

One invocation acts on source and target banks indexed by triples.
{[}O{]}, Section 2.1, uses central coordinates \(c_1,\ldots,c_h,c_*\),
with

\[
(Gx)_j=\sum_{T\ni j}x_T,\qquad (Gx)_*=\sum_Tx_T,
\quad (Rc)_S=\frac{\sum_{j\in S}c_j-c_*}{2}.
\]

But every source triple contributes to exactly three of the first \(h\)
sums. Thus \textbf{the source-dependent part} of \(c_*\) is redundant.
This is not a claim that arbitrary initial scratch obeys a linear
relation.

Use only \(h\) central coordinates, and define

\[
(G'x)_j=\sum_{T\ni j}x_T,
\qquad
(R'c)_S=\frac12\sum_{j\in S}c_j-\frac16\sum_{j=1}^h c_j.
\]

For each source triple \(T\) and target triple \(S\),

\[
(R'G')_{S,T}=\frac{|S\cap T|}{2}-\frac36
=\frac{|S\cap T|-1}{2}.
\]

This is exactly the old central matrix. The arbitrary-scratch argument
is also exact:

\[
-R'c+R'(c+G'x)=R'G'x,
\qquad(c+G'x)-G'x=c.
\]

No initially-zero or constrained center is required.

Keep the side corrections \(+1/2\) for intersection zero and \(-1/2\)
for intersection two. Their sum with the central coefficient is zero for
intersection sizes \(0,1,2\), and one for intersection size \(3\). The
scalar invocation remains the bank shear, and the three invocations
remain the signed bank exchange.

All surviving central gates retain their full local-space labels in
{[}O{]}, Table (2.7). The sole central decrease still has dimension
\(h\) per central wire, but there are now \(h\) such wires rather than
\(h+1\). Therefore the total loss is

\[
\boxed{L=Ih^2=3v^2h^2.}
\]

The denser expression for \(R'\) increases a fixed number of scalar
operations. Since \(h\) and the entire network are fixed before the
input length is chosen, these operations enter the additive linear-work
constant, not the recursion exponent. Its coefficients are explicit
rationals with nonzero denominators.

\section{3. The shared-sum replacement and its reversible role
compiler}\label{the-shared-sum-replacement-and-its-reversible-role-compiler}

Let

\[
(H_2x)_T=\sum_{S:\,|S\cap T|=2}x_S.
\]

For each fixed pair \(Q\), list the \(r=h-2\) remaining vertices, and
abbreviate \(x_i=x_{Q\cup\{i\}}\). Compute prefixes up to index \(r-3\)
and suffixes from index \(2\):

\[
P_i=x_0+\cdots+x_i\ (0\le i\le r-3),
\qquad S_i=x_i+\cdots+x_{r-1}\ (2\le i\le r-1).
\]

Inject two separate partial outputs into target \(t\):

\begin{itemize}
\tightlist
\item
  \(t=0\): \(x_1\) and \(S_2\);
\item
  \(1\le t\le r-2\): \(P_{t-1}\) and \(S_{t+1}\);
\item
  \(t=r-1\): \(P_{r-3}\) and \(x_{r-2}\).
\end{itemize}

Every required ordered source--target pair occurs exactly once, in its
unique common pair group. Every partial support has at most \(h-4\)
source triples. Do not combine the two partial outputs into one scratch
output.

The number of addition nodes and output uses is

\[
c=\binom h2(2h-10),\qquad q=2\binom h2(h-2).
\]

Assign a provisional physical role to every directed use, including
output uses. At each addition, identify its first incoming role with its
first outgoing role. The compiler first fans out sources, then processes
additions topologically: add the second incoming role into the first,
then copy that pivot into its other outgoing roles. Copies mean
reversible additions into arbitrary scratch, not overwrites. This
defines an invertible linear map \(A\) on the scratch roles. The dirty
contents will be canceled below.

There are \(2c+q\) directed uses and \(c\) pivot identifications, so

\[
\boxed{R_2=c+q=\binom h2(4h-14).}
\]

Let \(V_2\) inject each source into its pivot and let \(J\) read the
designated partial-output roles. Exact coefficient propagation gives
\(JAV_2=H_2\). Write \(J_2=-J/2\).

The two-pass computation

\[
Az\ \longrightarrow\ \text{subtract }J_2Az\ \longrightarrow\ z,
\]

followed by injection \(z\leftarrow z+V_2x\), a second application of
\(A\), positive output injection, reversal, and removal of \(V_2x\),
restores arbitrary scratch and adds

\[
J_2A(z+V_2x)-J_2Az=-\tfrac12H_2x.
\]

This proves the scalar identity. The next two sections address the
additional phase-label obligation; the scalar identity alone would not
suffice.

\section{4. Binary spaces and a complete forward invocation
schedule}\label{binary-spaces-and-a-complete-forward-invocation-schedule}

Use the local binary space \(D=\mathbb F_2^h\). Retain {[}O{]}'s stage
factors

\[
E=B\perp P,\qquad \dim P=1,\qquad\dim B=h^{j-1}-1,
\]

and its future norm-one line \(Q_f\). Suppress tensoring every displayed
label with \(Q_f\). For a local subspace \(U\), write

\[
\ell(U)=B\otimes U,\qquad
M(U)=(B\otimes D)\perp(P\otimes U),\qquad
H=E\otimes D=M(D).
\]

For the pair family \(Q=\{a,b\}\), its triple labels are
\(t_i=e_a+e_b+e_i\). They are orthonormal. For a partial support \(S\),
let \(U_S=\operatorname{span}\{t_i:i\in S\}\). The label attached to a
reversible addition or copy is the support of its underlying DAG node,
including both predecessors of an addition.

Along each physical role during the forward compiler, these supports
only grow. Each output role has one designated readout and its last
support is precisely the support it injects. The tagged compiler checks
these properties on the complete multi-pair graph, including fan-out of
a source into different pair groups.

Let \(z\) denote new shared scratch and \(a_0\) the unchanged
disjointness side roles. \(V_0,J_0\) are the original disjointness
injection and signed readout maps. The following sequence realizes one
\textbf{forward} invocation. A label in a bank-specific row uses that
bank entry's triple \(t\).

\begin{longtable}[]{@{}rll@{}}
\toprule\noalign{}
Step & Scalar operation & Common label on touched roles \\
\midrule\noalign{}
\endhead
\bottomrule\noalign{}
\endlastfoot
1 & \(y\leftarrow y-J_0 a_0\) & \(\ell(t_Y)\) \\
2 & \(z\leftarrow A z\) & \(\ell(D)\) on each elementary gate \\
3 & \(y\leftarrow y-J_2 z\) & \(\ell(D)\) \\
4 & \(z\leftarrow A^{-1} z\) & \(\ell(D)\) \\
5 & \(y\leftarrow y-R' c\) & \(\ell(D)\) \\
6 & \(a_0\leftarrow a_0+V_0 x\), \(z\leftarrow z+V_2 x\) & \(M(t_X)\) \\
7 & \(z\leftarrow A z\) & \(M(U_S)\) for each tagged gate \\
8 & \(c\leftarrow c+G' x\) & \(H\) \\
9 & \(y\leftarrow y+R' c\) & \(\ell(D)\) \\
10 & \(y\leftarrow y+J_0 a_0\), \(y\leftarrow y+J_2 z\) &
\(M(t_Y^\perp)\) \\
11 & \(c\leftarrow c-G' x\) & \(H\) \\
12 & \(z\leftarrow A^{-1} z\) & \(H\) \\
13 & \(a_0\leftarrow a_0-V_0 x\), \(z\leftarrow z-V_2 x\) & \(H\) \\
\end{longtable}

The disjointness negative readout is deliberately performed before the
new negative readout raises the Y-bank labels to \(\ell(D)\).

The input bank labels are \((E\otimes\langle t\rangle)\) on X and
\(\ell(t)\) on Y. They finish at \(H\) on X and \(M(t^\perp)\) on Y,
exactly the boundaries in {[}O{]}. All scratch begins at label zero and
finishes at \(H\) before its unchanged extension through the
future-space complement to the global sink label.

Every new-scratch edge is increasing. In particular, support propagation
in step 7 gives \(M(U_A)\subseteq M(U_S)\) when \(A\subseteq S\). An
output supported on \(S\) contains only sources orthogonal to its target
\(t\), so \(U_S\subseteq t^\perp\), permitting its step-10 readout. Step
12 raises any remaining label to \(H\) before uncomputation.

\section{5. The inverse invocation: do not reverse labels
blindly}\label{the-inverse-invocation-do-not-reverse-labels-blindly}

Stage 2 uses the inverse scalar invocation with logical source Y and
target X. Reversing the scalar updates is necessary; reversing the
forward label sequence would introduce decreases and is \textbf{not}
what is done.

The following chronological schedule is exactly the inverse scalar word
after exchanging logical source and target:

\begin{longtable}[]{@{}rll@{}}
\toprule\noalign{}
Step & Scalar operation & Common label \\
\midrule\noalign{}
\endhead
\bottomrule\noalign{}
\endlastfoot
1 & \(a_0\leftarrow a_0+V_0 y\), \(z\leftarrow z+V_2 y\) &
\(\ell(t_Y)\) \\
2 & \(z\leftarrow A z\) & \(\ell(U_S)\) for each tagged gate \\
3 & \(c\leftarrow c+G' y\) & \(\ell(D)\) \\
4 & \(x\leftarrow x-J_0 a_0\), \(x\leftarrow x-J_2 z\) & \(M(t_X)\) \\
5 & \(x\leftarrow x-R' c\) & \(H\) \\
6 & \(c\leftarrow c-G' y\) & \(\ell(D)\) \\
7 & \(z\leftarrow A^{-1} z\) & \(M(U_S^\perp)\) in reverse tagged-gate
order \\
8 & \(a_0\leftarrow a_0-V_0 y\), \(z\leftarrow z-V_2 y\) &
\(M(t_Y^\perp)\) \\
9 & \(x\leftarrow x+R' c\) & \(H\) \\
10 & \(z\leftarrow A z\) & \(H\) \\
11 & \(x\leftarrow x+J_2 z\) & \(H\) \\
12 & \(z\leftarrow A^{-1} z\) & \(H\) \\
13 & \(x\leftarrow x+J_0 a_0\) & \(H\) \\
\end{longtable}

A growing support path \(U_A\subseteq U_S\) becomes an increasing
complementary path \(U_S^\perp\subseteq U_A^\perp\) under reverse order
in step 7. At a designated output, the target \(t_X\) is orthogonal to
its source support \(U_S\), so \(M(t_X)\subseteq M(U_S^\perp)\). At a
source pivot the final transition to \(M(t_Y^\perp)\) is also
increasing.

Thus the inverse invocation preserves the required physical X/Y boundary
labels and adds no shared-scratch decreases. Its center alone drops from
\(H\) to \(\ell(D)\) at steps 5 to 6, once per central wire, of
dimension \(h\).

\section{6. Nondegeneracy, residual orthonormal bases, and the full
interface}\label{nondegeneracy-residual-orthonormal-bases-and-the-full-interface}

The previous inclusions must also have the residual bases required by
{[}O{]}, Lemma 2.2. Here is a classification, rather than an assumption
that nondegeneracy alone is enough in characteristic two.

\textbf{Subset and complement paths.} Every \(U_S\) is spanned by an
orthonormal subset. Residuals for \(U_A\subseteq U_S\) and
\(U_S^\perp\subseteq U_A^\perp\) have the explicit orthonormal basis
indexed by \(S\setminus A\).

\textbf{Complements of partial supports.} A partial support has at most
\(h-4\) members of one pair family. Its union of coordinate supports
misses at least two local coordinate units. Therefore \(U_S^\perp\) is
nondegenerate and contains a norm-one coordinate vector.

\textbf{Readout residuals.} The target label is a member of the same
orthonormal pair family, outside \(S\). The space

\[
(U_S\perp\langle t\rangle)^\perp
\]

is nondegenerate and contains a coordinate unit outside both the partial
support and the target: at least one such coordinate remains.
Consequently it is nonalternating and has an orthonormal basis by the
constructive unit-line/alternating-plane argument in {[}O{]}, Lemma 2.3.

\textbf{Crossings between lower and upper levels.} For example,

\[
\ell(U_A)\subset M(U_B^\perp)
\]

has residual

\[
(B\otimes U_A^\perp)\perp(P\otimes U_B^\perp).
\]

Both local complements are of the preceding kind. The analogous crossing
to \(M(t)\) has residual
\((B\otimes U_A^\perp)\perp(P\otimes\langle t\rangle)\). All nonzero
summands have norm-one vectors and orthonormal bases. Zero-dimensional
factors are omitted.

The relevant nonzero \(B\) spaces and future complements retain the
upstream norm-one-coordinate argument: their excluded tensor indicators
have support \(3^r<h^r\). Tensoring orthonormal bases and taking
orthogonal sums preserves an orthonormal basis. This includes the final
auxiliary extension to the full ambient \(h^3\)-dimensional space.

The old disjointness-side edges are exactly the upstream edges, whose
two triple supports occupy at most six coordinates; \(h=24\) is more
than enough. Every surviving central edge has the original label and
residual type. The only downward edges are therefore still the central
crossings, now on \(h\) wires per invocation.

The source/sink labels and signed scalar exchange remain unchanged. The
exceptional terminal vector has weight \(3^3=27\), so the same endpoint
translations in {[}O{]}, Proposition 2.4, give \(C^{\otimes mf}\) on
\textbf{every} physical role, including all new scratch roles. Each
scalar gate has a common frame on all roles it touches, so the
cancellation identity {[}O{]}, Equation (2.6), applies without assuming
anything about scratch values.

This establishes the proposed modified finite-network interface: \(s\)
directional steps, each implementable as copies of \(C^{\otimes f}\),
and a fixed number of pointwise operations and address permutations,
uniformly for every integer \(f\ge1\) and spectator bits.

\section{7. Exact network counts}\label{exact-network-counts}

For the reduced center plus shared-sum network,

\[
\begin{aligned}
W&=2v^3+3v^2\left(v\binom{h-3}{3}+\binom h2(4h-14)+h\right),\\
L&=3v^2h^2,\\
\Delta&=2(N-L)=2v^2(v-3h^2),\\
s&=Wm-\Delta.
\end{aligned}
\]

The first term in \(W\) counts the two banks; the next terms count
disjointness roles, shared intersection-two roles, and central roles in
all \(3v^2\) invocations.

At \(h=24\),

\[
\begin{aligned}
v&=2024,&m&=13824,&N&=8291469824,\\
I&=12289728,&L&=7078883328,\\
W&=33377983614976,&\Delta&=2425172992,\\
s&=461417243068255232.
\end{aligned}
\]

In particular \(\Delta>0\), and \(1<s/W<m\).

The \(h=24\) choice was selected by a finite parameter exploration, not
by a proof of global optimality. No optimality claim is needed for the
asserted upper bound.

\section{8. A no-padding tensor
recurrence}\label{a-no-padding-tensor-recurrence}

{[}O{]}, Section 2.6, pads \(W\) to a power of two to make every
recursive batch full. Reducing \(W\) alone need not improve its
exponent: both old and new widths may pad to the same power of two. This
is why the earlier local width reduction was not automatically an
exponent improvement.

\textbf{Batching lemma.} Suppose a fixed finite network with \(W\) roles
and \(s\) directional steps satisfies the complete interface above for a
fixed \(m\ge2\), with \(1<s/W<m\). Then its tensor transform has
deterministic charged cost

\[
O\!\left(2^k(k+1)^\theta\right),\qquad \theta=\log_m(s/W),
\]

without requiring that \(W\) be a power of two.

\textbf{Proof.} First transform \(W\) arbitrary arrays at once. For
\(k\ge m\), put \(f=\lfloor k/m\rfloor\) and \(r=k-mf\). Process the
fewer than \(m\) leftover axes directly. One directional step contains
\(F=2^{k-f}\) independent fibers of length \(2^f\). After the upstream
linear-time permutation, process \(\lfloor F/W\rfloor\) full batches
recursively. Process the remaining \(F\bmod W<W\) fibers with the
ordinary coordinate-by-coordinate tensor algorithm, at cost at most

\[
B(W-1)f2^f
\]

for an absolute \(B\). The quotient, remainder, batch boundaries, and
the ordinary leftover transforms use the same logarithmic-word model;
discovery and movement remain linear in array length.

Let \(T(k)\) be the cost for \(W\) arrays and \(t(k)=T(k)/(W2^k)\). With
all nonrecursive fixed-network work charged,

\[
T(k)\le A_0W2^k+s\left\lfloor\frac{2^{k-f}}W\right\rfloor T(f)
+B s(W-1)f2^f.
\]

Using the upper bound on the floor gives

\[
t(k)\le A_0+\frac{s}{W}t(f)+B s f2^{f-k}.
\]

Since \(k\ge mf\) and \(m\ge2\),

\[
f2^{f-k}\le f2^{-(m-1)f}\le1.
\]

Thus the leftover work is a bounded additive term in this normalized
recurrence:

\[
\boxed{t(k)\le A+\frac{s}{W}t(\lfloor k/m\rfloor).}
\]

Its constants are fixed independently of \(k\). The finite range \(k<m\)
is handled directly. The same recurrence unrolling as in {[}O{]},
Theorem 2.6, gives the critical exponent \(\theta=\log_m(s/W)\). The
recursive arrays can be processed serially with workspace reuse, and
their address fields require \(O(k+1)\) bits. A single requested
transform initializes the other \(W-1\) roles once, at cost
\(O(W2^k)=O(2^k)\) because \(W\) is fixed. This proves the lemma.

This is a deterministic remainder-handling argument, not an assumption
of free padding, fractional batches, or average-case input data.

\section{9. Propagation to all Fourier
lengths}\label{propagation-to-all-fourier-lengths}

The tensor interface is exactly the input consumed by {[}O{]},
Proposition 4.2. Its local Fourier compiler in Section 3 uses the same
fixed two-coordinate complex matrix \(C\); none of its local-width,
scalar-preparation, nonzero-denominator, or root requirements change.

The sector packing work is unchanged. It gives cost

\[
O\!\left(R(\ell+1)^\theta(1+\log r_{\max})^4\right)
+\operatorname{poly}(\ell,r_{\max},\log R).
\]

The small-prime working length, CRT permutations, and charged chirp
convolution in Section 5 are likewise unchanged. Substituting
\(\ell=\Theta(\log n/\log\log n)\) gives

\[
\boxed{T(n)=O\!\left(n(\log n)^\theta(\log\log n)^{4-\theta}\right).}
\]

This also accounts for the transformed fixed convolution operand. The
specified root of order less than \(1024 n^3\) remains sufficient. New
fixed rational coefficients introduce no additional root requirement.

Write

\[
a=1-\theta=-\frac{\log(1-\Delta/(mW))}{\log m}.
\]

For \textbf{every fixed} \(0<\delta<a\), the \(\log\log\) factor is
absorbed by \((\log n)^{a-\delta}\), yielding
\(O(n(\log n)^{1-\delta})\). There is no mandatory factor-of-two loss in
\(\delta\); the upstream headline's halving was a convenient
conservative corollary.

\section{10. A rational certificate for the proposed
exponent}\label{a-rational-certificate-for-the-proposed-exponent}

The new critical saving is

\[
a=5.5127258849324290013\ldots\times10^{-10}.
\]

A safe witness is \(\delta=55/10^{11}=5.5\times10^{-10}\). This is
certified without a floating-point comparison.

Set \(\epsilon=\Delta/(mW)\). The verifier establishes by exact
fractions that

\[
\epsilon>\frac{55}{10^{11}}\frac{477}{50},\qquad \log(13824)<\frac{477}{50}=9.54.
\]

For the latter inequality, the positive finite sum

\[
\sum_{j=0}^{40}\frac{(477/50)^j}{j!}>13824
\]

certifies \(\exp(9.54)>13824\). Therefore

\[
a=\frac{-\log(1-\epsilon)}{\log m}
>\frac{\epsilon}{9.54}>\frac{55}{10^{11}}.
\]

The reported finer numerical values are computed using exact rational
atanh-series enclosures, with an explicit tail bound. The JSON contains
the actual rational interval endpoints. Decimal endpoints can round to
the same displayed string; the rational endpoints remain distinct
rigorous enclosures.

\clearpage
\section{11. Separate the improvements and their proof
dependencies}\label{separate-the-improvements-and-their-proof-dependencies}

Here \(a=1-\theta\) is the critical saving \textbf{before} absorbing the
\(\log\log\) factor. Any smaller positive \(\delta\) gives a pure-power
corollary.

\begin{longtable}[]{@{}
  >{\raggedright\arraybackslash}p{(\columnwidth - 4\tabcolsep) * \real{0.3000}}
  >{\raggedleft\arraybackslash}p{(\columnwidth - 4\tabcolsep) * \real{0.4000}}
  >{\raggedright\arraybackslash}p{(\columnwidth - 4\tabcolsep) * \real{0.3000}}@{}}
\toprule\noalign{}
\begin{minipage}[b]{\linewidth}\raggedright
Construction
\end{minipage} & \begin{minipage}[b]{\linewidth}\raggedleft
Critical saving \(a\)
\end{minipage} & \begin{minipage}[b]{\linewidth}\raggedright
Additional proof ingredients
\end{minipage} \\
\midrule\noalign{}
\endhead
\bottomrule\noalign{}
\endlastfoot
Uploaded paper, \(h=100\), padded & \(2.10643843004018\times10^{-13}\) &
Upstream \\
Original network, \(h=24\), padded & \(3.95761002385911\times10^{-10}\)
& Generalize the displayed fixed-size network to this valid \(h\) \\
Reduced center, \(h=24\), padded & \(5.22969896327091\times10^{-10}\) &
Section 2 of this note; upstream batching retained \\
Reduced center, \(h=24\), unpadded & \(5.30775793289321\times10^{-10}\)
& Also the remainder-handling lemma \\
Reduced center + shared sums, \(h=24\), unpadded &
\(5.51272588493243\times10^{-10}\) & Also the two-orientation phase
schedules \\
\end{longtable}

The first reduced-center version supports the simpler
\(\delta=5.2\times10^{-10}\) claim without relying on the proposed
shared-sum phase integration. The next supports \(5.3\times10^{-10}\).
Thus a flaw in the shared-sum argument would not, by itself, revoke the
simpler candidates.

Most of the gain relative to the uploaded paper comes from a better
ground size. That must not be portrayed as a new fundamental Fourier
method. The shared-sum component adds about \textbf{3.8617\%} to the
reduced-center unpadded critical saving. Relative to the already-retuned
\(h=25\) complex-network threshold of
\(4.18479903721200\times10^{-10}\), the new critical saving is only
about \textbf{31.73\% larger}, not thousands of times larger. The direct
all-length transfer and the new finite modifications must be
distinguished from the ground-size retuning already present in Colkitt's
work. All ratios concern asymptotic exponent savings, \textbf{not
measured speedups}.

\section{12. Reproduction, evidence, and remaining
review}\label{reproduction-evidence-and-remaining-review}

Run with Python 3.10 or later, using only its standard library:

\begin{Shaded}
\begin{Highlighting}[]
\ExtensionTok{python}\NormalTok{ verify\_extension.py}
\end{Highlighting}
\end{Shaded}

The program performs:

\begin{enumerate}
\def\labelenumi{\arabic{enumi}.}
\item
  Exact reconstruction of the uploaded \(h=100\) counts, exact new
  counts, and rational witness comparisons.
\item
  The reduced-center coefficient identity for all intersection sizes,
  and the general dirty-scratch cancellation argument recorded in its
  output.
\item
  Exact verification of every local \(H_2\) coefficient at \(h=24\) and
  \(h=25\), plus dirty-scratch executions and agreement of tagged and
  untagged compiler implementations.
\item
  A support-label trace for every role of each full multi-pair graph,
  including source fan-out and unique designated readout roles.
\item
  Actual binary linear algebra on every canonical-pair wire edge in the
  forward first/third stages and inverse middle stage at both ground
  sizes. It verifies inclusions, nondegeneracy, residual dimensions, and
  orthonormal bases for the relevant nonalternating components.
  Coordinate permutations cover the other pair groups; the global
  support trace checks the shared source fan-out between groups.
\item
  Execution of the complete modified local scalar word, its reversed
  middle word, and their three-shear signed bank exchange on arbitrary
  rational scratch examples at \(h=7,8\). These small sizes test the
  scalar identities only, not positive asymptotic savings.
\end{enumerate}

At \(h=24\), the shared graph has \(22,632\) roles, \(31,096\)
elementary operations in its forward scratch map, and \(12,144\) output
reads. Its local \(H_2\) matrix has \(2,024^2=4,096,576\) entries;
\(127,512\) are nonzero. At \(h=25\) the prior corresponding counts are
also reproduced.

\textbf{Limits of these checks.} They are not a numerical execution of
the astronomical complete DFT algorithm. They do not mechanically prove
the asymptotic recurrence, the upstream exact-width compiler, the entire
global \(h^3\)-dimensional phase network, or the all-length reduction.
Those conclusions use the written arguments and identified upstream
results. The canonical frame audit is a symmetry-reduced check of the
modified component, not a claim to enumerate every global role of the
full network.

Before presenting this publicly as an established advance, independent
review should prioritize the two orientation schedules, the transfer
from full support traces to global residual classification, the
arbitrary-scratch interpretation, and the no-padding recurrence. A
comparison with other follow-ups is still needed for priority. The
formulas in Sections 2 and 8 are short enough to review separately from
the larger shared-sum construction.

\section{References}\label{references}

{[}O{]} OpenAI. \emph{An explicit power saving for the exact discrete
Fourier transform}. September 25, 2026. Uploaded file
\texttt{main\ (1).pdf}. Especially Sections 2.1--2.6, Proposition 3.1,
Proposition 4.2, and Section 5.4.

{[}F{]} OpenAI. \emph{Finite tensor savings and exact Fourier circuits}.
September 25, 2026. Uploaded file \texttt{main.pdf}. The independent
finite-win route is not a dependency of the proposed explicit
quantitative extension.

{[}D{]} Douglas Colkitt. \texttt{CrocSwap/integer-mult-bounds},
\texttt{notes/independent-complex.tex} and
\texttt{docs/research/current-status.md}, accessed October 8, 2026. The
latter states that his results remain conditional and not
independently/formally verified.

https://github.com/CrocSwap/integer-mult-bounds/blob/main/notes/independent-complex.tex

https://github.com/CrocSwap/integer-mult-bounds/blob/main/docs/research/current-status.md

{[}P{]} The October 7 local shared-sum working note and executable
verifiers supplied earlier in this conversation. They explicitly did not
establish the Fourier transfer. The present note supersedes that
uncertainty with the proposed complete schedules and transfer argument;
it does not retrospectively strengthen what the earlier checks
established.

\end{document}
